% Generated draft: requires mathematical and source review.
\documentclass[11pt]{article}
\usepackage[margin=25mm]{geometry}
\usepackage{fontspec}
\setmainfont{Latin Modern Roman}
\usepackage{amsmath,amssymb,mathtools,unicode-math}
\setmathfont{Latin Modern Math}
\usepackage{enumitem,needspace,xurl,hyperref}
\setlength{\parindent}{0pt}
\setlength{\parskip}{5pt}
\setlength{\emergencystretch}{3em}
\newcommand{\N}{\mathbb{N}}
\newcommand{\Z}{\mathbb{Z}}
\newcommand{\Q}{\mathbb{Q}}
\newcommand{\R}{\mathbb{R}}
\newcommand{\C}{\mathbb{C}}
\newcommand{\F}{\mathbb{F}}
\newcommand{\A}{\mathbb{A}}
\newcommand{\PP}{\mathbb{P}}
\newcommand{\E}{\mathbb{E}}
\newcommand{\eps}{\varepsilon}
\newcommand{\dd}{\,\mathrm{d}}
\begin{document}
\section*{UnsolvedMath ID: 701713}
\textbf{Source problem code:} catalogue entry
\subsection{Definitions and mathematical statement}
For a group \(K\), define its derived series by \(K^{(0)}=K\) and
\[
K^{(i+1)}=[K^{(i)},K^{(i)}],
\]
where \([L,L]\) is the subgroup generated by all \(x^{-1}y^{-1}xy\) with \(x,y\in L\). The group \(K\) is soluble if \(K^{(d)}=\{1\}\) for some nonnegative integer \(d\); its derived length \(\operatorname{dl}(K)\) is the least such \(d\).

Interpret ``character table'' as the ordinary complex character table. Write \(\operatorname{Irr}(K)\) for the set of irreducible complex characters of a finite group \(K\), and \(\operatorname{Cl}(K)\) for its set of conjugacy classes. For \(C\in\operatorname{Cl}(K)\), let \(\chi(C)\) denote the common value of \(\chi\) on \(C\). Two finite groups \(G,H\) have identical character tables if there are bijections
\[
\alpha:\operatorname{Irr}(G)\longrightarrow\operatorname{Irr}(H),\qquad
\beta:\operatorname{Cl}(G)\longrightarrow\operatorname{Cl}(H)
\]
such that
\[
\chi(C)=\alpha(\chi)\bigl(\beta(C)\bigr)
\quad\text{for every }\chi\in\operatorname{Irr}(G)\text{ and }C\in\operatorname{Cl}(G).
\]
The supplied record does not explicitly specify the character field or the convention for table identity; the interpretation above is the ordinary complex convention, allowing row and column permutations.

\textbf{Question (S. Mattarei).} Do there exist finite soluble groups \(G\) and \(H\) with identical character tables and
\[
\operatorname{dl}(G)=2,\qquad \operatorname{dl}(H)=4?
\]
Equivalently, the required derived-series conditions are \(G^{(2)}=\{1\}\), \(G^{(1)}\ne\{1\}\), \(H^{(4)}=\{1\}\), and \(H^{(3)}\ne\{1\}\).

The supplied background reports that, for every integer \(n\ge2\), pairs of nonisomorphic finite soluble groups with identical character tables and derived lengths \(n\) and \(n+1\) were constructed by S. Mattarei in the cited 1995 paper. This reported construction concerns consecutive derived lengths and is distinct from the requested pair of lengths \(2\) and \(4\).
\subsection{Short English statement}
Do there exist two finite soluble groups with the same ordinary complex character table, up to permutations of rows and columns, whose derived lengths are respectively \(2\) and \(4\)? The supplied background reports examples with consecutive lengths \(n,n+1\) for every integer \(n\ge2\).
\subsection{Sources}
\begin{itemize}
\item Local UnsolvedMath registry, permanent ID \(701713\), exact source problem code \texttt{catalogue entry}, title ``Kourovka Notebook Problem 20.55.'' The supplied record attributes the problem to S. Mattarei, gives the proposed year as 2022, and labels its status as open. This status is supplied registry data, not an independently verified assessment.
\item \emph{The Kourovka Notebook (Unsolved Problems in Group Theory)}, problem 20.55. The supplied background identifies the 20th issue (2022) and the 21st issue edition (2026), with identifier \texttt{arXiv:1401.0300v48}. The supplied collection description names E. I. Khukhro and V. D. Mazurov as editors. These bibliographic details have not been independently verified.
\item S. Mattarei, \emph{J. Algebra} \textbf{175} (1995), 157--178. The supplied statement cites this paper for constructions with derived lengths \(n\) and \(n+1\) for every integer \(n\ge2\). No article title or direct link is supplied.
\end{itemize}
\end{document}
